\documentclass[a4paper, 11pt]{article}

% === 폰트 ===
\usepackage{fontspec}
\usepackage[hangul]{luatexko}

% Latin 폰트 (숫자, 영문, 기호)
\setmainfont[Ligatures=TeX]{TeX Gyre Heros}
\setsansfont{TeX Gyre Heros}
\setmonofont{Noto Sans Mono}

% 한글 폰트
\setmainhangulfont{Noto Sans CJK KR}[
  BoldFont={Noto Sans CJK KR Bold}
]
\setsanshangulfont{Noto Sans CJK KR}[
  BoldFont={Noto Sans CJK KR Bold}
]

% === 레이아웃 ===
\usepackage{geometry}
\geometry{left=25mm, right=25mm, top=30mm, bottom=30mm, headheight=15pt}
\usepackage{graphicx}
\usepackage{xcolor}
\usepackage{tcolorbox}
\usepackage{booktabs}
\usepackage{array}
\usepackage{longtable}
\usepackage{colortbl}
\usepackage{calc}
\usepackage{fancyhdr}
\usepackage{titlesec}
\usepackage{hyperref}
\usepackage{emoji}
\usepackage{amsmath}
\usepackage{amssymb}
\usepackage{enumitem}
\usepackage{multirow}
\usepackage{makecell}
\usepackage{float}

% === 색상 ===
\definecolor{mainblue}{HTML}{1976D2}
\definecolor{lightblue}{HTML}{E3F2FD}
\definecolor{lightgray}{HTML}{F5F5F5}
\definecolor{warnbg}{HTML}{FFF8E1}

% === 이모지 ===
\setemojifont{Noto Color Emoji}

% === 헤더/풋터 ===
\pagestyle{fancy}
\fancyhf{}
\fancyhead[L]{\small AI디지털리터러시 — 1주차 2차시}
\fancyhead[R]{\small 강의노트 v1.3.0.0}
\fancyfoot[C]{\thepage}
\renewcommand{\headrulewidth}{0.4pt}
\renewcommand{\footrulewidth}{0pt}

% === 섹션 스타일 ===
\titleformat{\section}
  {\Large\bfseries\color{mainblue}}
  {\thesection.}{0.5em}{}[\titlerule]
\titleformat{\subsection}
  {\large\bfseries}{}{0em}{}
\titleformat{\subsubsection}
  {\normalsize\bfseries}{}{0em}{}

% === tcolorbox ===
\tcbuselibrary{skins, breakable}
\newtcolorbox{infobox}[1][]{
  colback=lightblue, colframe=mainblue,
  fonttitle=\bfseries, title=#1,
  breakable, enhanced}
\newtcolorbox{tipbox}[1][]{
  colback=lightgray, colframe=gray,
  fonttitle=\bfseries, title=#1,
  breakable, enhanced}
\newtcolorbox{warnbox}[1][]{
  colback=warnbg, colframe=orange,
  fonttitle=\bfseries, title=#1,
  breakable, enhanced}

% === hyperref ===
\hypersetup{
  unicode=true,
  pdfencoding=unicode,
  colorlinks=true,
  linkcolor=mainblue,
  urlcolor=mainblue
}

% === pandoc 호환 ===
\providecommand{\tightlist}{%
  \setlength{\itemsep}{0pt}\setlength{\parskip}{0pt}}

% blockquote 스타일
\renewenvironment{quote}
  {\begin{tcolorbox}[colback=lightblue!30, colframe=mainblue!50,
   left=8pt, right=8pt, top=4pt, bottom=4pt, breakable]}
  {\end{tcolorbox}}

% code block 스타일

\begin{document}

% === 제목 ===
\begin{center}
{\LARGE\bfseries\color{mainblue} AI디지털리터러시 — 1주차 2차시 강의노트}\\[0.5em]
{\Large\bfseries 기계학습의 원리}\\[0.5em]
{\normalsize 문서 버전: v1.3.0.0}
\end{center}

\vspace{1em}
\hrule
\vspace{1em}

\hypertarget{uxc218uxc5c5-uxc815uxbcf4}{%
\subsection{수업 정보}\label{uxc218uxc5c5-uxc815uxbcf4}}

\begin{longtable}{p{3cm}p{11cm}}
\toprule\noalign{}
항목 & 내용 \\
\midrule\noalign{}
\endhead
\bottomrule\noalign{}
\endlastfoot
\textbf{교과목명} & AI디지털리터러시 \\
\textbf{주차/차시} & 1주차 2차시 \\
\textbf{주제} & 기계학습의 원리 \\
\textbf{수업 시간} & 50분 (1시간) \\
\textbf{수업 대상} & 청주교육대학교 4학년 \\
\textbf{CK 강화 목표} & 기계학습의 원리(지도/비지도/강화학습)를
이해한다. \\
\textbf{관련 성취기준} & {[}6실05-05{]} \\
\end{longtable}

\hypertarget{uxc218uxc5c5-uxd750uxb984-uxc694uxc57d}{%
\subsection{수업 흐름
요약}\label{uxc218uxc5c5-uxd750uxb984-uxc694uxc57d}}

\begin{longtable}{p{4.5cm}p{4.5cm}p{5.0cm}}
\toprule\noalign{}
\textbf{시간} & \textbf{단계} & \textbf{활동 내용} \\
\midrule\noalign{}
\endhead
\bottomrule\noalign{}
\endlastfoot
5분 & \emoji{clapper-board} 도입 & AI, ML, DL의 포함 관계 \\
15분 & \emoji{blue-book} 전개 1 & 기계학습의 정의: 데이터로부터 패턴을 학습하는 AI
기술 \\
15분 & \emoji{green-book} 전개 2 & 학습 유형: 지도학습(분류/예측), 비지도학습(군집화),
강화학습(보상) \\
10분 & \emoji{orange-book} 전개 3 & 머신러닝 학습 과정: 데이터 수집 → 학습 → 평가 →
예측 \\
5분 & \emoji{memo} 정리 & 지도학습 = 엔트리 AI 모델의 원리 연결, 다음 주 퀴즈
안내 \\
\end{longtable}

\hypertarget{uxb3c4uxc785-ai-ml-dluxc758-uxd3ecuxd568-uxad00uxacc4-5uxbd84}{%
\section{\emoji{clapper-board} 도입: AI, ML, DL의 포함 관계
(5분)}\label{uxb3c4uxc785-ai-ml-dluxc758-uxd3ecuxd568-uxad00uxacc4-5uxbd84}}

\hypertarget{uxcc28uxc2dc-uxbcf5uxc2b5}{%
\subsection{1차시 복습}\label{uxcc28uxc2dc-uxbcf5uxc2b5}}

\begin{quote}
\emoji{pushpin} \textbf{핵심 질문}: ``1차시에서 배운 AI, 그 안에는 무엇이 있을까요?''
\end{quote}

지난 시간에 우리는 AI의 정의와 유형을 배웠어요. AI는 ``인간의 지능을
모방하는 컴퓨터 시스템''이라고 했죠. 그런데 AI를 구현하는 방법은 여러
가지가 있어요. 오늘은 그중에서도 가장 중요한 \textbf{기계학습(Machine
Learning)}에 대해 알아보겠습니다.

\hypertarget{ai-ml-dluxc758-uxd3ecuxd568-uxad00uxacc4}{%
\subsection{AI, ML, DL의 포함
관계}\label{ai-ml-dluxc758-uxd3ecuxd568-uxad00uxacc4}}

\textbf{{[}그림 01{]} AI, ML, DL의 포함 관계}

AI, ML, DL은 자주 혼동되는 용어입니다. 이 세 가지의 관계를 정확히
이해하는 것이 중요해요.

\begin{verbatim}
AI (인공지능)
 ⊃ ML (기계학습)
    ⊃ DL (딥러닝)
\end{verbatim}

\hypertarget{ai-artificial-intelligence-uxc778uxacf5uxc9c0uxb2a5}{%
\subsubsection{\emoji{robot} AI (Artificial Intelligence,
인공지능)}\label{ai-artificial-intelligence-uxc778uxacf5uxc9c0uxb2a5}}

\begin{itemize}
\tightlist
\item
  \textbf{가장 큰 개념}: 인간의 지능을 모방하는 모든 기술
\item
  \textbf{범위}: 규칙 기반 시스템부터 최신 딥러닝까지 모두 포함
\item
  \textbf{예시}:

  \begin{itemize}
  \tightlist
  \item
    규칙 기반: 체스 프로그램 (if-then 규칙)
  \item
    기계학습: 스팸 필터, 추천 시스템
  \item
    전문가 시스템: 1980년대 의료 진단 시스템
  \end{itemize}
\end{itemize}

\hypertarget{ml-machine-learning-uxae30uxacc4uxd559uxc2b5}{%
\subsubsection{\emoji{brain} ML (Machine Learning,
기계학습)}\label{ml-machine-learning-uxae30uxacc4uxd559uxc2b5}}

\begin{itemize}
\tightlist
\item
  \textbf{AI의 하위 개념}: AI를 구현하는 대표적 방법
\item
  \textbf{핵심 특징}: 데이터로부터 패턴을 \textbf{자동으로 학습}
\item
  \textbf{차이점}: 사람이 규칙을 일일이 작성하지 않음
\item
  \textbf{예시}:

  \begin{itemize}
  \tightlist
  \item
    이메일 스팸 필터
  \item
    유튜브 추천 알고리즘
  \item
    날씨 예측 모델
  \end{itemize}
\end{itemize}

\hypertarget{dl-deep-learning-uxb525uxb7ecuxb2dd}{%
\subsubsection{\emoji{fire} DL (Deep Learning,
딥러닝)}\label{dl-deep-learning-uxb525uxb7ecuxb2dd}}

\begin{itemize}
\tightlist
\item
  \textbf{ML의 하위 개념}: 기계학습의 한 종류
\item
  \textbf{핵심 특징}: 인공신경망을 여러 층으로 쌓아 복잡한 패턴 학습
\item
  \textbf{``Deep''의 의미}: 신경망의 층(Layer)이 깊다(많다)
\item
  \textbf{예시}:

  \begin{itemize}
  \tightlist
  \item
    ChatGPT (자연어 처리)
  \item
    AlphaGo (바둑)
  \item
    자율주행 (이미지 인식)
  \end{itemize}
\end{itemize}

\begin{quote}
\emoji{light-bulb} \textbf{핵심 통찰}: AI ⊃ ML ⊃ DL 관계를 이해하면, 각 기술의 범위와
특성을 명확히 구분할 수 있습니다.
\end{quote}

\hypertarget{uxbe44uxc720uxb85c-uxc774uxd574uxd558uxae30}{%
\subsubsection{비유로
이해하기}\label{uxbe44uxc720uxb85c-uxc774uxd574uxd558uxae30}}

AI를 ``차량''에 비유하면: - \textbf{AI} = 모든 차량 (자전거, 오토바이,
자동차, 비행기) - \textbf{ML} = 엔진이 있는 차량 (오토바이, 자동차,
비행기) - \textbf{DL} = 첨단 전기차 (테슬라)

기계학습은 AI의 일부이고, 딥러닝은 기계학습의 일부예요.

\hypertarget{uxc804uxac1c-1-uxae30uxacc4uxd559uxc2b5uxc758-uxc815uxc758-15uxbd84}{%
\section{\emoji{blue-book} 전개 1: 기계학습의 정의
(15분)}\label{uxc804uxac1c-1-uxae30uxacc4uxd559uxc2b5uxc758-uxc815uxc758-15uxbd84}}

\hypertarget{uxae30uxacc4uxd559uxc2b5uxc774uxb780}{%
\subsection{기계학습이란?}\label{uxae30uxacc4uxd559uxc2b5uxc774uxb780}}

\begin{quote}
\textbf{기계학습(Machine Learning)}: 컴퓨터가 데이터로부터 패턴을 스스로
학습하는 AI 기술
\end{quote}

기계학습은 1959년 아서 새뮤얼(Arthur Samuel)이 정의한 개념입니다:
``명시적으로 프로그래밍하지 않고도 컴퓨터가 학습할 수 있는 능력''

\hypertarget{uxc804uxd1b5uxc801-uxd504uxb85cuxadf8uxb798uxbc0d-vs-uxae30uxacc4uxd559uxc2b5}{%
\subsection{전통적 프로그래밍 vs
기계학습}\label{uxc804uxd1b5uxc801-uxd504uxb85cuxadf8uxb798uxbc0d-vs-uxae30uxacc4uxd559uxc2b5}}

기계학습을 이해하려면, 먼저 전통적 프로그래밍과 어떻게 다른지 알아야
해요.

\hypertarget{uxc804uxd1b5uxc801-uxd504uxb85cuxadf8uxb798uxbc0d}{%
\subsubsection{전통적
프로그래밍}\label{uxc804uxd1b5uxc801-uxd504uxb85cuxadf8uxb798uxbc0d}}

\begin{verbatim}
규칙 + 데이터 → 결과
\end{verbatim}

\textbf{예시: 스팸 메일 필터 (전통적 방식)}

\begin{enumerate}
\def\labelenumi{\arabic{enumi}.}
\tightlist
\item
  프로그래머가 규칙 작성:

  \begin{itemize}
  \tightlist
  \item
    IF ``무료'' 포함 AND ``당첨'' 포함 → 스팸
  \item
    IF ``대출'' 포함 AND ``즉시'' 포함 → 스팸
  \item
    IF 발신자가 블랙리스트 → 스팸
  \end{itemize}
\item
  문제점:

  \begin{itemize}
  \tightlist
  \item
    규칙을 모두 나열하기 불가능
  \item
    새로운 스팸 기법에 대응 어려움
  \item
    유지보수가 복잡
  \end{itemize}
\end{enumerate}

\hypertarget{uxae30uxacc4uxd559uxc2b5}{%
\subsubsection{기계학습}\label{uxae30uxacc4uxd559uxc2b5}}

\begin{verbatim}
데이터 + 결과 → 규칙 (패턴)
\end{verbatim}

\textbf{예시: 스팸 메일 필터 (기계학습 방식)}

\begin{enumerate}
\def\labelenumi{\arabic{enumi}.}
\tightlist
\item
  컴퓨터에게 수천 개의 메일 제공:

  \begin{itemize}
  \tightlist
  \item
    스팸 메일 5,000개 (정답: 스팸)
  \item
    정상 메일 5,000개 (정답: 정상)
  \end{itemize}
\item
  컴퓨터가 스스로 패턴 발견:

  \begin{itemize}
  \tightlist
  \item
    ``무료'', ``당첨''이 자주 나오면 스팸일 확률 높음
  \item
    친구 이메일에는 이름이 자주 나옴
  \item
    스팸은 링크가 많음
  \end{itemize}
\item
  장점:

  \begin{itemize}
  \tightlist
  \item
    사람이 규칙을 작성하지 않아도 됨
  \item
    새로운 스팸 기법도 자동으로 학습
  \item
    데이터가 많을수록 성능 향상
  \end{itemize}
\end{enumerate}

\hypertarget{uxc65c-uxae30uxacc4uxd559uxc2b5uxc774-uxd544uxc694uxd55cuxac00}{%
\subsection{왜 기계학습이
필요한가?}\label{uxc65c-uxae30uxacc4uxd559uxc2b5uxc774-uxd544uxc694uxd55cuxac00}}

기계학습이 필요한 이유는 크게 세 가지예요.

\hypertarget{uxaddcuxce59uxc744-uxc791uxc131uxd558uxae30-uxc5b4uxb824uxc6b4-uxbb38uxc81c}{%
\subsubsection{\textbf{①} 규칙을 작성하기 어려운
문제}\label{uxaddcuxce59uxc744-uxc791uxc131uxd558uxae30-uxc5b4uxb824uxc6b4-uxbb38uxc81c}}

\textbf{예시: 고양이 vs 개 구분}

전통적 프로그래밍으로 한다면: - IF 귀가 뾰족하면 → 고양이? - 하지만
치와와도 귀가 뾰족함 - IF 수염이 있으면 → 고양이? - 하지만 일부 개도
수염이 있음

→ 규칙을 명확히 작성하기 불가능!

기계학습으로 한다면: - 고양이 사진 1,000장 + 개 사진 1,000장 제공 -
컴퓨터가 스스로 패턴 발견 - 우리가 모르는 미세한 차이까지 학습

\hypertarget{uxbcf5uxc7a1uxd55c-uxd328uxd134uxc774-uxc788uxb294-uxbb38uxc81c}{%
\subsubsection{\textbf{②} 복잡한 패턴이 있는
문제}\label{uxbcf5uxc7a1uxd55c-uxd328uxd134uxc774-uxc788uxb294-uxbb38uxc81c}}

\textbf{예시: 음성 인식}

\begin{itemize}
\tightlist
\item
  같은 단어도 사람마다 발음이 다름
\item
  억양, 속도, 음높이가 천차만별
\item
  규칙으로 모든 경우를 다루기 불가능
\item
  기계학습으로 수백만 개 음성 데이터 학습
\end{itemize}

\hypertarget{uxd658uxacbduxc774-uxacc4uxc18d-uxbcc0uxd558uxb294-uxbb38uxc81c}{%
\subsubsection{\textbf{③} 환경이 계속 변하는
문제}\label{uxd658uxacbduxc774-uxacc4uxc18d-uxbcc0uxd558uxb294-uxbb38uxc81c}}

\textbf{예시: 주식 가격 예측}

\begin{itemize}
\tightlist
\item
  경제 상황이 계속 변함
\item
  과거 규칙이 미래에 안 맞을 수 있음
\item
  기계학습은 새 데이터로 계속 업데이트
\end{itemize}

\hypertarget{uxc77cuxc0c1-uxc18d-uxae30uxacc4uxd559uxc2b5-uxc608uxc2dc}{%
\subsection{일상 속 기계학습
예시}\label{uxc77cuxc0c1-uxc18d-uxae30uxacc4uxd559uxc2b5-uxc608uxc2dc}}

우리는 이미 기계학습을 매일 사용하고 있어요.

\begin{longtable}{p{4cm}p{10cm}}
\toprule\noalign{}
서비스 & 기계학습 활용 \\
\midrule\noalign{}
\endhead
\bottomrule\noalign{}
\endlastfoot
\emoji{e-mail} \textbf{Gmail 스팸 필터} & 수십억 개의 메일 데이터로 스팸 패턴
학습 \\
\emoji{television} \textbf{유튜브 추천} & 시청 기록으로 취향 분석, 관련 영상 추천 \\
\emoji{speaking-head} \textbf{Siri, Alexa} & 수백만 명의 음성으로 음성 인식 학습 \\
\emoji{camera-with-flash} \textbf{얼굴 인식} & 수천 장의 얼굴 사진으로 특징 학습 \\
\emoji{automobile} \textbf{자율주행} & 수백만 킬로미터 주행 데이터로 운전 학습 \\
\emoji{shopping-cart} \textbf{쿠팡 추천} & 구매 기록으로 선호 상품 예측 \\
\end{longtable}

\begin{quote}
\emoji{light-bulb} \textbf{핵심 메시지}: 기계학습은 ``사람이 규칙을 작성하는 대신,
컴퓨터가 데이터에서 규칙을 찾는'' 기술입니다.
\end{quote}

\hypertarget{uxc804uxac1c-2-uxae30uxacc4uxd559uxc2b5uxc758-3uxac00uxc9c0-uxd559uxc2b5-uxc720uxd615-15uxbd84}{%
\section{\emoji{green-book} 전개 2: 기계학습의 3가지 학습 유형
(15분)}\label{uxc804uxac1c-2-uxae30uxacc4uxd559uxc2b5uxc758-3uxac00uxc9c0-uxd559uxc2b5-uxc720uxd615-15uxbd84}}

기계학습은 크게 세 가지 방식으로 학습합니다. 각각의 특징을 이해해
봅시다.

\hypertarget{uxc9c0uxb3c4uxd559uxc2b5-supervised-learning}{%
\subsection{\textbf{①} 지도학습 (Supervised
Learning)}\label{uxc9c0uxb3c4uxd559uxc2b5-supervised-learning}}

\begin{quote}
\textbf{지도학습}: 정답이 있는 데이터로 학습하는 방식
\end{quote}

\hypertarget{uxd575uxc2ec-uxac1cuxb150}{%
\subsubsection{핵심 개념}\label{uxd575uxc2ec-uxac1cuxb150}}

``선생님이 문제와 정답을 주고 공부하는 것''과 같아요.

\textbf{학습 과정}: 1. 데이터 + 정답(Label) 제공 2. 컴퓨터가 데이터와
정답의 관계 학습 3. 새로운 데이터가 들어오면 정답 예측

\hypertarget{uxc9c0uxb3c4uxd559uxc2b5uxc758-uxb450-uxac00uxc9c0-uxc885uxb958}{%
\subsubsection{지도학습의 두 가지
종류}\label{uxc9c0uxb3c4uxd559uxc2b5uxc758-uxb450-uxac00uxc9c0-uxc885uxb958}}

\hypertarget{uxbd84uxb958-classification}{%
\paragraph{\emoji{bar-chart} 분류 (Classification)}\label{uxbd84uxb958-classification}}

\textbf{입력 → 범주(카테고리)}

\begin{longtable}{p{3cm}p{4cm}p{7cm}}
\toprule\noalign{}
예시 & 입력 & 출력 (범주) \\
\midrule\noalign{}
\endhead
\bottomrule\noalign{}
\endlastfoot
이메일 스팸 판별 & 메일 내용 & 스팸 / 정상 \\
질병 진단 & 증상, 검사 결과 & 독감 / 감기 / 정상 \\
꽃 종류 분류 & 꽃잎 길이, 너비 등 & 장미 / 튤립 / 백합 \\
고양이 vs 개 & 사진 & 고양이 / 개 \\
\end{longtable}

\hypertarget{uxc608uxce21-regression---uxd68cuxadc0}{%
\paragraph{\emoji{chart-increasing} 예측 (Regression) -
회귀}\label{uxc608uxce21-regression---uxd68cuxadc0}}

\textbf{입력 → 숫자(연속적 값)}

\emoji{light-bulb} \textbf{쉬운 이해}: - \textbf{회귀}는 ``\textbf{얼마나?}''라는 질문에
답하는 것 - 연속적인 숫자를 맞히는 학습

\textbf{교수 멘트}:

\begin{verbatim}
"회귀(Regression)라는 말이 어렵게 느껴지죠?
간단하게 말하면 '숫자를 맞히는 것'이에요.

예를 들어볼까요?
- 이 집은 **얼마**일까? → 3억 5천만원
- 시험 **몇 점** 받을까? → 85점
- 내일 기온이 **몇 도**일까? → 18도

이렇게 숫자로 된 답을 맞히는 게 회귀예요!"
\end{verbatim}

\textbf{회귀 vs 분류 비교표}:

\begin{longtable}{p{2cm}p{4cm}p{8cm}}
\toprule\noalign{}
구분 & 회귀 & 분류 \\
\midrule\noalign{}
\endhead
\bottomrule\noalign{}
\endlastfoot
\textbf{질문} & ``얼마나?'' & ``어느 것?'' \\
\textbf{답 형태} & 숫자 (연속적) & 범주 (구분) \\
\textbf{예시 1} & 시험 점수: 85점 & 합격/불합격: 합격 \\
\textbf{예시 2} & 집값: 3억 5천 & 아파트/빌라: 아파트 \\
\textbf{예시 3} & 기온: 18도 & 더위/추위: 더위 \\
\textbf{예시 4} & 과일 무게: 180g & 과일 종류: 사과 \\
\end{longtable}

\textbf{강조}: - 회귀는 0\textasciitilde100점,
0원\textasciitilde10억원처럼 \textbf{연속적인 범위}의 숫자 - 분류는
합격/불합격, 고양이/개처럼 \textbf{정해진 범주} 중 하나

\begin{longtable}{p{3cm}p{4cm}p{7cm}}
\toprule\noalign{}
예시 & 입력 & 출력 (숫자) \\
\midrule\noalign{}
\endhead
\bottomrule\noalign{}
\endlastfoot
집값 예측 & 평수, 위치, 연식 & 5억 2천만 원 \\
성적 예측 & 공부 시간, 출석률 & 85점 \\
기온 예측 & 시간, 계절, 습도 & 23.5도 \\
주식 가격 예측 & 과거 가격, 거래량 & 52,300원 \\
\end{longtable}

\hypertarget{uxc9c0uxb3c4uxd559uxc2b5-uxc608uxc2dc-uxace0uxc591uxc774-vs-uxac1c}{%
\subsubsection{지도학습 예시: 고양이 vs
개}\label{uxc9c0uxb3c4uxd559uxc2b5-uxc608uxc2dc-uxace0uxc591uxc774-vs-uxac1c}}

\textbf{데이터 준비}: - 고양이 사진 1,000장 (정답: 고양이) - 개 사진
1,000장 (정답: 개)

\textbf{AI가 보는 데이터의 모습} (특징 표 예시):

\begin{longtable}{p{1.5cm}p{2cm}p{2.5cm}p{2.5cm}p{2.5cm}p{3cm}}
\toprule\noalign{}
사진 번호 & 귀 모양 & 수염 길이 & 눈 크기 & 털 색깔 & 라벨 (정답) \\
\midrule\noalign{}
\endhead
\bottomrule\noalign{}
\endlastfoot
1 & 뾰족 & 길다 & 크다 & 갈색 & 고양이 \emoji{check-mark-button} \\
2 & 접힘 & 짧다 & 작다 & 검정 & 개 \emoji{check-mark-button} \\
3 & 뾰족 & 길다 & 크다 & 흰색 & 고양이 \emoji{check-mark-button} \\
4 & 접힘 & 짧다 & 중간 & 갈색 & 개 \emoji{check-mark-button} \\
\ldots{} & \ldots{} & \ldots{} & \ldots{} & \ldots{} & \ldots{} \\
\end{longtable}

\emoji{light-bulb} \textbf{AI가 발견하는 패턴}:

\begin{verbatim}
"귀 모양이 뾰족하고 수염이 길면 → 고양이!"
"귀가 접히고 수염이 짧으면 → 개!"
\end{verbatim}

\textbf{학습 과정}: 1. 컴퓨터에게 사진 + 정답 제공 2. 컴퓨터가 고양이와
개의 \textbf{특징(Feature)} 학습 - 특징(Feature) = 데이터의 속성 (귀
모양, 눈 모양, 털 패턴 등) 3. 특징들의 패턴을 스스로 발견

\textbf{예측}: - 새로운 동물 사진 입력 - 특징 추출 (귀 모양, 수염 길이
등) - 학습한 패턴과 비교 - 컴퓨터가 ``고양이'' 또는 ``개'' 판별

\hypertarget{uxc5d4uxd2b8uxb9ac-ai-uxbaa8uxb378-uxc5f0uxacb0}{%
\subsubsection{엔트리 AI 모델
연결}\label{uxc5d4uxd2b8uxb9ac-ai-uxbaa8uxb378-uxc5f0uxacb0}}

\emoji{light-bulb} \textbf{중요}: 3주차에 배울 \textbf{엔트리 AI 이미지 분류 모델}이
바로 지도학습입니다!

엔트리에서 여러분이 직접: 1. 고양이/개 사진 찍기 (데이터 수집) 2. ``이건
고양이'', ``이건 개'' 라벨 붙이기 (정답 지정) 3. ``모델 학습하기'' 버튼
클릭 (학습) 4. 새 사진으로 테스트 (예측)

→ 이 전체 과정이 지도학습이에요!

\hypertarget{uxbe44uxc9c0uxb3c4uxd559uxc2b5-unsupervised-learning}{%
\subsection{\textbf{②} 비지도학습 (Unsupervised
Learning)}\label{uxbe44uxc9c0uxb3c4uxd559uxc2b5-unsupervised-learning}}

\begin{quote}
\textbf{비지도학습}: 라벨(정답 표시) 없이 데이터의 패턴을 스스로 찾는
방식
\end{quote}

\hypertarget{uxc815uxb2f5uxc774-uxc5c6uxb2e4uxc758-uxc815uxd655uxd55c-uxc758uxbbf8}{%
\subsubsection{``정답이 없다''의 정확한
의미}\label{uxc815uxb2f5uxc774-uxc5c6uxb2e4uxc758-uxc815uxd655uxd55c-uxc758uxbbf8}}

\emoji{light-bulb} \textbf{매우 중요한 개념!}

\textbf{교수 전략}:

\begin{verbatim}
"비지도학습을 설명할 때 '정답이 없다'고 하는데,
이게 정확히 무슨 뜻일까요?

여러분, 잘 들어보세요.
'정답이 없다' = '라벨(정답 스티커)이 없다'는 뜻이에요!"
\end{verbatim}

\hypertarget{uxc9c0uxb3c4uxd559uxc2b5-vs-uxbe44uxc9c0uxb3c4uxd559uxc2b5-uxb370uxc774uxd130-uxad6cuxc870-uxbe44uxad50}{%
\subsubsection{지도학습 vs 비지도학습 데이터 구조
비교}\label{uxc9c0uxb3c4uxd559uxc2b5-vs-uxbe44uxc9c0uxb3c4uxd559uxc2b5-uxb370uxc774uxd130-uxad6cuxc870-uxbe44uxad50}}

\textbf{판서 또는 슬라이드}:

\begin{verbatim}
지도학습 (Supervised Learning):
  데이터 + 라벨(정답)

  [고양이 사진] + "고양이" \emoji{check-mark-button}
  [개 사진] + "개" \emoji{check-mark-button}
  
  → 모든 데이터에 정답 스티커가 붙어 있음!

비지도학습 (Unsupervised Learning):
  데이터만 (라벨 없음)

  고객A: 25세, 구매액 10만원, 패션 선호
  고객B: 45세, 구매액 200만원, 전자제품 선호
  고객C: 28세, 구매액 15만원, 패션 선호
  
  → 어떤 그룹인지 정답 스티커가 없음! \emoji{cross-mark}
  → AI가 스스로 비슷한 사람끼리 묶어서 그룹 발견!
\end{verbatim}

\textbf{강조할 점}:

\begin{verbatim}
지도학습: "이건 고양이"라는 라벨이 있음
비지도학습: "이 고객은 알뜰형"이라는 라벨이 없음

하지만 라벨이 없어도 AI는 스스로 패턴을 찾아
의미 있는 그룹을 만들어요!
\end{verbatim}

\hypertarget{uxc65c-uxb77cuxbca8uxc774-uxc5c6uxc744uxae4c-3uxac00uxc9c0-uxc774uxc720}{%
\subsubsection{왜 라벨이 없을까? (3가지
이유)}\label{uxc65c-uxb77cuxbca8uxc774-uxc5c6uxc744uxae4c-3uxac00uxc9c0-uxc774uxc720}}

\textbf{교수 전략}: 각 이유를 구체적 예시와 함께 설명

\hypertarget{uxb77cuxbca8-uxc790uxccb4uxac00-uxc5c6uxc744-uxb54c}{%
\paragraph{\textbf{①} 라벨 자체가 없을
때}\label{uxb77cuxbca8-uxc790uxccb4uxac00-uxc5c6uxc744-uxb54c}}

\textbf{예시: 쇼핑몰 신규 고객}

\begin{verbatim}
상황:
  쇼핑몰에 신규 고객 10,000명이 가입했어요.

문제:
  이 고객들이 어떤 취향인지 아직 몰라요.
  
  "이 고객은 알뜰형"
  "이 고객은 명품 선호형"
  
  이런 라벨을 붙일 수가 없어요!
  왜? → 아직 정보가 없으니까!

해결:
  \emoji{cross-mark} 지도학습 불가능: 라벨이 없으니까
  \emoji{check-mark-button} 비지도학습: 구매 패턴만 보고 자동 그룹화
\end{verbatim}

\textbf{교수 멘트}:

\begin{verbatim}
"여러분, 신규 고객의 취향을 어떻게 알 수 있을까요?
구매 이력이 없는데 '이 사람은 이런 스타일'이라고
라벨을 붙일 수가 없잖아요!

하지만 AI는 나이, 첫 구매 품목, 검색 키워드 등을 보고
비슷한 사람끼리 묶어서 그룹을 만들어요."
\end{verbatim}

\hypertarget{uxb77cuxbca8-uxb9ccuxb4e4uxae30-uxc5b4uxb824uxc6b8-uxb54c}{%
\paragraph{\textbf{②} 라벨 만들기 어려울
때}\label{uxb77cuxbca8-uxb9ccuxb4e4uxae30-uxc5b4uxb824uxc6b8-uxb54c}}

\textbf{예시: 뉴스 기사 분류}

\begin{verbatim}
상황:
  인터넷에 매일 수백만 개의 뉴스 기사가 올라와요.

문제:
  각 기사에 "정치", "경제", "스포츠" 라벨을 붙이려면?
  
  → 사람이 일일이 읽고 분류해야 함
  → 하루에 100만 개? 불가능!
  → 시간/비용이 너무 많이 듦

해결:
  \emoji{cross-mark} 지도학습: 라벨 만들기 불가능 (시간/비용)
  \emoji{check-mark-button} 비지도학습: 단어 패턴만 보고 자동 분류
\end{verbatim}

\textbf{교수 멘트}:

\begin{verbatim}
"여러분이 매일 올라오는 100만 개 뉴스에
일일이 '정치', '경제', '스포츠' 라벨을 붙인다고 상상해보세요.

하루 종일 해도 못 끝나겠죠?

하지만 AI는 기사에서 '국회', '의원' 같은 단어가 많으면
정치 그룹으로, '주식', '환율' 같은 단어가 많으면
경제 그룹으로 자동으로 묶어요!"
\end{verbatim}

\hypertarget{uxc228uxaca8uxc9c4-uxd328uxd134uxc744-uxbc1cuxacacuxd558uxace0-uxc2f6uxc744-uxb54c}{%
\paragraph{\textbf{③} 숨겨진 패턴을 발견하고 싶을
때}\label{uxc228uxaca8uxc9c4-uxd328uxd134uxc744-uxbc1cuxacacuxd558uxace0-uxc2f6uxc744-uxb54c}}

\textbf{예시: 새로운 고객 그룹 발견}

\begin{verbatim}
상황:
  기존에 알려진 고객 유형:
  - 10대 학생
  - 20대 직장인
  - 30대 부모

문제:
  하지만 이것만으로는 부족해요!
  사람도 몰랐던 새로운 패턴이 있을 수 있어요.

AI 발견:
  - "오전 출근 시간에 편의점에서 샐러드 사는 그룹"
    → 건강 관심형 직장인
  
  - "주말 저녁에 온라인 쇼핑하는 그룹"
    → 여가 시간 쇼핑형
  
  - "새벽 3시에 라면 구매하는 그룹"
    → 야간 근무자 or 수험생

해결:
  사람이 미리 정의하지 않은 새로운 패턴을
  AI가 데이터에서 스스로 발견!
\end{verbatim}

\textbf{교수 멘트}:

\begin{verbatim}
"이게 비지도학습의 진짜 힘이에요!

사람이 생각하지 못한 패턴을 AI가 발견해요.

예를 들어 넷플릭스는 비지도학습으로
'금요일 밤 10시에 로맨스 영화 보는 사람들'
이런 새로운 시청 패턴을 발견했어요.

이건 사람이 미리 라벨로 만들 수 없는 거예요!"
\end{verbatim}

\hypertarget{uxad70uxc9d1uxd654-clustering-uxbe44uxc9c0uxb3c4uxd559uxc2b5uxc758-uxb300uxd45c-uxbc29uxbc95}{%
\subsubsection{군집화 (Clustering): 비지도학습의 대표
방법}\label{uxad70uxc9d1uxd654-clustering-uxbe44uxc9c0uxb3c4uxd559uxc2b5uxc758-uxb300uxd45c-uxbc29uxbc95}}

\textbf{개념 설명}:

\begin{verbatim}
군집화 = 유사한 데이터끼리 묶어서 그룹(클러스터) 만들기
\end{verbatim}

\textbf{3단계 과정}:

\begin{verbatim}
1단계: 입력
  고객 데이터만 제공 (라벨 없이)
  - 고객A: 나이 25, 구매액 10만원, 패션 선호
  - 고객B: 나이 45, 구매액 200만원, 전자제품 선호
  - 고객C: 나이 28, 구매액 15만원, 패션 선호
  ...

2단계: 과정
  AI가 유사한 고객끼리 자동 그룹화
  (라벨 없이 패턴만으로!)
  
  어떻게?
  - 나이가 비슷한 사람
  - 구매 금액이 비슷한 사람
  - 선호 카테고리가 비슷한 사람
  → 이런 기준으로 자동 분류

3단계: 출력
  그룹 A, B, C... (AI가 스스로 발견한 그룹)
  
  사람이 나중에 이름 붙임:
  - 그룹 A → "20대 패션 선호형"
  - 그룹 B → "40대 전자제품 선호형"
  - 그룹 C → "30대 생활용품 선호형"
\end{verbatim}

\hypertarget{uxad6cuxccb4uxc801-uxc608uxc2dc-3uxac00uxc9c0}{%
\subsubsection{구체적 예시
3가지}\label{uxad6cuxccb4uxc801-uxc608uxc2dc-3uxac00uxc9c0}}

\hypertarget{uxc608uxc2dc-1-uxc1fcuxd551uxbab0-uxace0uxac1d-uxc138uxbd84uxd654-uxc0c1uxc138}{%
\paragraph{예시 1: 쇼핑몰 고객 세분화
(상세)}\label{uxc608uxc2dc-1-uxc1fcuxd551uxbab0-uxace0uxac1d-uxc138uxbd84uxd654-uxc0c1uxc138}}

\begin{verbatim}
┌─────────────────────────────────────┐
│ 예시 1: 쇼핑몰 고객 세분화          │
└─────────────────────────────────────┘

상황:
  신규 고객 10,000명
  취향 라벨 없음 (어떤 고객인지 모름)

\emoji{cross-mark} 지도학습이 불가능한 이유:
  "이 고객은 알뜰형"이라는 라벨이 없음!
  
\emoji{check-mark-button} 비지도학습 적용:
  1. 데이터 입력: 나이, 구매액, 선호 카테고리
  2. AI 분석: 구매 패턴 자동 분석
  3. 자동 그룹화: 3~5개 그룹 발견

결과:
  그룹 1 (30%): 알뜰형
    - 특징: 저가 제품, 할인 쿠폰 활용
    - 나이: 20~30대
    - 구매액: 평균 5만원
  
  그룹 2 (20%): 명품 선호형
    - 특징: 고가 제품, 브랜드 중시
    - 나이: 30~50대
    - 구매액: 평균 50만원
  
  그룹 3 (25%): 트렌드 민감형
    - 특징: 신상품, 유행 제품
    - 나이: 10~20대
    - 구매액: 평균 15만원
  
  그룹 4 (25%): 실속형
    - 특징: 가성비 중시, 리뷰 확인
    - 나이: 전 연령
    - 구매액: 평균 10만원

활용:
  각 그룹에 맞춤 마케팅!
  - 알뜰형 → 할인 쿠폰 발송
  - 명품 선호형 → 신상 명품 소개
  - 트렌드 민감형 → 최신 트렌드 정보
  - 실속형 → 가성비 제품 추천
\end{verbatim}

\textbf{교수 팁}:

\begin{verbatim}
"이 예시를 설명할 때:

1. 먼저 '신규 고객은 라벨이 없다'를 강조
2. 지도학습이 왜 불가능한지 명확히
3. 비지도학습으로 어떻게 해결하는지 설명
4. 결과가 실제 비즈니스에 어떻게 활용되는지 연결

학생들이 '아, 이래서 비지도학습이 필요하구나!'
하고 납득할 수 있도록!"
\end{verbatim}

\hypertarget{uxc608uxc2dc-2-uxb137uxd50cuxb9aduxc2a4-uxcd94uxcc9c-uxc2dcuxc2a4uxd15c}{%
\paragraph{예시 2: 넷플릭스 추천
시스템}\label{uxc608uxc2dc-2-uxb137uxd50cuxb9aduxc2a4-uxcd94uxcc9c-uxc2dcuxc2a4uxd15c}}

\begin{verbatim}
┌─────────────────────────────────────┐
│ 예시 2: 넷플릭스 추천               │
└─────────────────────────────────────┘

상황:
  "좋아요" 데이터만으로는 취향 파악 어려움
  사용자가 무엇을 좋아할지 정확한 라벨 없음

비지도학습 적용:
  1. 시청 기록 수집
     - 어떤 영화를 봤나?
     - 언제 봤나?
     - 얼마나 봤나? (중간에 껐나?)
  
  2. 비슷한 시청 패턴을 가진 사용자끼리 그룹화
  
  3. 그룹 내 다른 사람들이 본 영화 추천

결과:
  "이 영화를 본 사람은 저 영화도 봤어요"
  
  → 이게 바로 비지도학습 결과!

왜 비지도학습?
  - "이 사람은 액션 선호형"이라는 명확한 라벨 없음
  - 사람의 취향은 복잡함 (액션+코미디+로맨스)
  - AI가 패턴을 발견해서 비슷한 사람 찾음
\end{verbatim}

\hypertarget{uxc608uxc2dc-3-uxb274uxc2a4-uxc790uxb3d9-uxbd84uxb958}{%
\paragraph{예시 3: 뉴스 자동
분류}\label{uxc608uxc2dc-3-uxb274uxc2a4-uxc790uxb3d9-uxbd84uxb958}}

\begin{verbatim}
┌─────────────────────────────────────┐
│ 예시 3: 뉴스 자동 분류              │
└─────────────────────────────────────┘

상황:
  매일 수백만 개 뉴스 기사 생성
  일일이 라벨 붙이기 불가능

비지도학습 적용:
  1. 뉴스 기사 텍스트 입력
  
  2. 단어 패턴 분석
     - "국회", "의원", "법안" → ?
     - "주식", "환율", "기업" → ?
     - "축구", "야구", "승리" → ?
  
  3. 유사한 단어 패턴끼리 자동 그룹화

결과:
  그룹 A: 정치 기사 (국회, 선거, 정당...)
  그룹 B: 경제 기사 (주식, 환율, 기업...)
  그룹 C: 스포츠 기사 (축구, 야구, 올림픽...)
  그룹 D: IT 기사 (AI, 스마트폰, 앱...)

활용:
  - 자동 카테고리 분류
  - 유사 기사 묶기
  - 뉴스 추천 시스템
\end{verbatim}

\hypertarget{uxcd08uxb4f1uxad50uxc721-uxbe44uxc720}{%
\subsubsection{초등교육
비유}\label{uxcd08uxb4f1uxad50uxc721-uxbe44uxc720}}

\textbf{교수 전략}:

\begin{verbatim}
"초등학생에게 어떻게 설명할까요?"
\end{verbatim}

\textbf{비유 1: 친구 그룹 나누기}

\begin{verbatim}
상황:
  교실에 30명의 친구가 있어요.
  
지도학습:
  선생님이 "철수는 활발한 그룹", "영희는 조용한 그룹"
  이렇게 정답표를 미리 만들어 줌

비지도학습:
  선생님의 정답표가 없어요!
  
  AI가 스스로 관찰:
  - 쉬는 시간에 운동장 가는 친구들
  - 도서관에서 책 읽는 친구들
  - 미술실에서 그림 그리는 친구들
  
  → 비슷한 친구들끼리 자동으로 그룹을 만들어요!
\end{verbatim}

\textbf{비유 2: 과일 분류}

\begin{verbatim}
지도학습:
  엄마가 "이건 사과", "이건 바나나"라고 알려주며 학습

비지도학습:
  과일을 보고 스스로:
  - 빨간색 과일끼리 묶기
  - 노란색 과일끼리 묶기
  - 둥근 과일끼리 묶기
  
  정답표 없이 패턴으로 분류!
\end{verbatim}

\hypertarget{uxc624uxac1cuxb150-uxb300uxc751}{%
\subsubsection{오개념 대응}\label{uxc624uxac1cuxb150-uxb300uxc751}}

\begin{longtable}{p{4.5cm}p{4.5cm}p{5.0cm}}
\toprule\noalign{}
\textbf{학생 오개념} & \textbf{올바른 이해} & \textbf{대응 방법} \\
\midrule\noalign{}
\endhead
\bottomrule\noalign{}
\endlastfoot
``비지도학습은 아무렇게나 해도 된다'' & 의미 있는 패턴을 찾아야 함 &
``라벨이 없어도 품질 기준은 있어요'' \\
``정답이 없으면 틀려도 괜찮다'' & 결과의 품질은 평가 가능 & ``실루엣
스코어 등으로 품질 측정'' \\
``지도학습보다 성능이 낮다'' & 목적이 다름 (발견 vs 예측) & ``새로운
패턴 발견에 강점'' \\
``비지도학습은 쓸모없다'' & 많은 실제 활용 사례 존재 & ``추천 시스템,
이상 탐지 등'' \\
\end{longtable}

\textbf{대응 전략}:

\begin{verbatim}
"여러분, 비지도학습은 '정답이 없어서 아무렇게나'가 아니에요!

오히려 '정답을 모를 때' 또는 '새로운 패턴을 찾고 싶을 때'
사용하는 강력한 도구예요!

넷플릭스, 유튜브, 쇼핑몰 추천이 모두 비지도학습을 사용해요.
정답 라벨이 없어도 의미 있는 결과를 만들어내죠!"
\end{verbatim}

\hypertarget{uxac15uxd654uxd559uxc2b5-reinforcement-learning}{%
\subsection{\textbf{③} 강화학습 (Reinforcement
Learning)}\label{uxac15uxd654uxd559uxc2b5-reinforcement-learning}}

\begin{quote}
\textbf{강화학습}: 시행착오를 통해 보상을 최대화하는 방향으로 학습
\end{quote}

\hypertarget{uxd575uxc2ec-uxac1cuxb150-1}{%
\subsubsection{핵심 개념}\label{uxd575uxc2ec-uxac1cuxb150-1}}

``게임하면서 점수로 배우기''

\textbf{학습 과정}: 1. AI가 행동(Action) 선택 2. 환경이 보상(Reward)
제공 - 좋은 행동 → + 보상 - 나쁜 행동 → - 보상 3. 보상을 최대화하는
방향으로 학습 4. 반복하며 최적의 전략 발견

\hypertarget{uxac15uxd654uxd559uxc2b5uxc758-uxad6cuxc131-uxc694uxc18c}{%
\subsubsection{강화학습의 구성
요소}\label{uxac15uxd654uxd559uxc2b5uxc758-uxad6cuxc131-uxc694uxc18c}}

\begin{longtable}{p{4.5cm}p{4.5cm}p{5.0cm}}
\toprule\noalign{}
\textbf{요소} & \textbf{설명} & \textbf{예시 (게임)} \\
\midrule\noalign{}
\endhead
\bottomrule\noalign{}
\endlastfoot
\textbf{Agent (에이전트)} & 학습하는 주체 & 게임 플레이어 \\
\textbf{Environment (환경)} & 에이전트가 상호작용하는 세계 & 게임
화면 \\
\textbf{State (상태)} & 현재 상황 & 점수, 위치, 체력 \\
\textbf{Action (행동)} & 에이전트가 선택할 수 있는 것 & 위/아래/좌/우
이동 \\
\textbf{Reward (보상)} & 행동에 대한 피드백 & 점수 +10, -5 \\
\end{longtable}

\hypertarget{uxc608uxc2dc-1-alphago-uxbc14uxb451}{%
\subsubsection{예시 1: AlphaGo
(바둑)}\label{uxc608uxc2dc-1-alphago-uxbc14uxb451}}

\textbf{목표}: 바둑에서 이기기

\textbf{학습 과정}: 1. AlphaGo가 바둑돌 놓기 (행동) 2. 게임 결과에 따라
보상: - 이기면 → +1 (큰 보상) - 지면 → -1 (큰 벌점) - 좋은 수 → +0.1
(작은 보상) 3. 수백만 번 자기 자신과 대국하며 학습 4. 이길 확률을 높이는
전략 발견

\textbf{결과}: 2016년 이세돌 9단을 4:1로 격파

\hypertarget{uxc608uxc2dc-2-uxc790uxc728uxc8fcuxd589}{%
\subsubsection{예시 2:
자율주행}\label{uxc608uxc2dc-2-uxc790uxc728uxc8fcuxd589}}

\textbf{목표}: 안전하고 효율적으로 운전

\textbf{학습 과정}: 1. AI가 핸들 조작 (행동) 2. 결과에 따라 보상: - 차선
유지 → +1 - 차선 이탈 → -5 - 사고 → -100 - 목적지 도착 → +10 3.
시뮬레이션에서 수백만 번 주행하며 학습 4. 보상을 최대화하는 운전 기술
습득

\hypertarget{uxc608uxc2dc-3-uxac8cuxc784-ai}{%
\subsubsection{예시 3: 게임 AI}\label{uxc608uxc2dc-3-uxac8cuxc784-ai}}

\textbf{게임}: 슈퍼마리오

\textbf{학습 과정}: 1. AI가 버튼 누르기 (점프, 이동) 2. 보상: - 코인
획득 → +10 - 적 제거 → +50 - 게임 오버 → -100 - 목표 도달 → +1000 3.
수천 번 게임하며 학습 4. 최고 점수를 얻는 플레이 방법 발견

\hypertarget{uxc2dcuxac01uxd654uxb85c-uxc774uxd574uxd558uxae30}{%
\subsubsection{시각화로
이해하기}\label{uxc2dcuxac01uxd654uxb85c-uxc774uxd574uxd558uxae30}}

\hypertarget{uxadf8uxb9bc-19-uxac15uxc544uxc9c0-uxd6c8uxb828uxc73cuxb85c-uxc774uxd574uxd558uxae30}{%
\paragraph{\texorpdfstring{\textbf{{[}그림 02{]} 강아지 훈련으로
이해하기}
\emoji{dog}}{{[}그림 02{]} 강아지 훈련으로 이해하기 \emoji{dog}}}\label{uxadf8uxb9bc-19-uxac15uxc544uxc9c0-uxd6c8uxb828uxc73cuxb85c-uxc774uxd574uxd558uxae30}}

강화학습의 원리를 일상 생활의 강아지 훈련으로 쉽게 설명합니다.

\textbf{교수 전략}:

\begin{verbatim}
"강화학습을 강아지 훈련으로 생각해볼까요?

1. 상황 (State): 강아지가 마당에 있어요
2. 행동 (Action): '앉아!' 명령 → 강아지가 앉음
3. 보상 (Reward): 잘했어! 간식 제공! \emoji{bone}
4. 학습 (Learning): '앉아 = 간식'으로 기억

다음번엔? 앉아 명령만 들어도 바로 앉아요!
왜? 간식을 받았던 기억 때문이죠.

AI도 똑같이 배워요!
- 좋은 결과 → 그 방법 기억
- 나쁜 결과 → 다른 방법 시도"
\end{verbatim}

\textbf{초등학생 설명}:

\begin{verbatim}
"여러분이 강아지를 훈련시킨다고 생각해보세요.
강아지가 앉으면 간식을 주고,
앉지 않으면 아무것도 안 주죠?

그럼 강아지는 어떻게 할까요?
간식을 받기 위해 계속 앉으려고 하겠죠!

AI도 똑같아요!
좋은 점수를 받으면 그 방법을 기억하고,
나쁜 점수를 받으면 다른 방법을 시도해요."
\end{verbatim}

\hypertarget{uxadf8uxb9bc-20-uxac8cuxc784-uxc810uxc218-uxd68duxb4dd-uxacfcuxc815}{%
\paragraph{\texorpdfstring{\textbf{{[}그림 03{]} 게임 점수 획득 과정}
\emoji{video-game}}{{[}그림 03{]} 게임 점수 획득 과정 \emoji{video-game}}}\label{uxadf8uxb9bc-20-uxac8cuxc784-uxc810uxc218-uxd68duxb4dd-uxacfcuxc815}}

게임 AI의 시행착오 학습 과정을 시각화합니다.

\textbf{교수 전략}:

\begin{verbatim}
"게임으로 더 재미있게 이해해볼까요?

게임 목표: 장애물 피하면서 최대한 멀리 가기

시도 1 (실패):
- 장애물 앞 → 점프 → 부딪힘! \emoji{collision}
- 점수: -10점 \emoji{cross-mark}
- 학습: "점프는 위험해..."

시도 2 (성공):
- 장애물 앞 → 옆으로 피함 → 성공! \emoji{sparkles}
- 점수: +50점 \emoji{check-mark-button}
- 학습: "옆으로 피하는 게 좋구나!"

100만 번 반복하면?
- 1,000번: 50점 \emoji{crying-face}
- 10,000번: 300점 \emoji{slightly-smiling-face}
- 100,000번: 850점 \emoji{smiling-face-with-smiling-eyes}
- 1,000,000번: 2,500점 \emoji{trophy}

→ 최고 점수 달성!

AI는 지치지 않고 24시간 연습할 수 있어요!"
\end{verbatim}

\textbf{오개념 방지}:

\begin{longtable}{p{5.0cm}p{9.0cm}}
\toprule\noalign{}
\textbf{학생 질문} & \textbf{교사 답변} \\
\midrule\noalign{}
\endhead
\bottomrule\noalign{}
\endlastfoot
``선생님, AI는 한 번에 잘하나요?'' & ``아니요! 처음엔 우리보다
못해요.하지만 백만 번 연습하면서 점점 좋아지죠.결국 우리보다 훨씬 잘하게
돼요!'' \\
``정답을 알려주나요?'' & ``아니요! 점수만 줘요.'이렇게 해라'가 아니라
'이건 50점이야'라고만 말해주죠.그럼 AI가 스스로 점수 높이는 방법을
찾아요!'' \\
``사람처럼 생각하나요?'' & ``아니요! 단순히 '높은 점수 = 좋은 행동'으로
배워요.의미를 이해하는 건 아니에요.'' \\
\end{longtable}

\hypertarget{uxbe44uxc720uxb85c-uxc774uxd574uxd558uxae30-1}{%
\subsubsection{비유로
이해하기}\label{uxbe44uxc720uxb85c-uxc774uxd574uxd558uxae30-1}}

\textbf{지도학습}: 선생님이 정답 알려주며 공부 \textbf{비지도학습}:
혼자서 패턴 찾기 \textbf{강화학습}: 게임하며 점수로 배우기 (시행착오)

\hypertarget{uxc138-uxac00uxc9c0-uxd559uxc2b5-uxc720uxd615-uxbe44uxad50}{%
\subsection{세 가지 학습 유형
비교}\label{uxc138-uxac00uxc9c0-uxd559uxc2b5-uxc720uxd615-uxbe44uxad50}}

\begin{longtable}{p{2.5cm}p{3.5cm}p{4.0cm}p{4.0cm}}
\toprule\noalign{}
\textbf{구분} & \textbf{지도학습} & \textbf{비지도학습} & \textbf{강화학습} \\
\midrule\noalign{}
\endhead
\bottomrule\noalign{}
\endlastfoot
\textbf{정답} & \emoji{hollow-red-circle} 있음 & \emoji{cross-mark} 없음 & \emoji{bullseye} 보상 \\
\textbf{목표} & 정답 맞추기 & 패턴 찾기 & 보상 최대화 \\
\textbf{비유} & 시험 공부 & 그룹 나누기 & 게임 연습 \\
\textbf{데이터} & 입력 + 정답 & 입력만 & 상태 + 행동 + 보상 \\
\textbf{예시} & 스팸 필터질병 진단집값 예측 & 고객 세분화뉴스 분류추천
시스템 & AlphaGo자율주행게임 AI \\
\end{longtable}

\begin{quote}
\emoji{light-bulb} \textbf{핵심 메시지}: 문제 유형에 따라 적절한 학습 방식을 선택하는
것이 중요합니다. 초등학교 AI 교육에서는 주로 \textbf{지도학습}(엔트리
이미지 분류)을 다룹니다.
\end{quote}

\hypertarget{uxc804uxac1c-3-uxba38uxc2e0uxb7ecuxb2dd-uxd559uxc2b5-uxacfcuxc815-10uxbd84}{%
\section{\emoji{orange-book} 전개 3: 머신러닝 학습 과정
(10분)}\label{uxc804uxac1c-3-uxba38uxc2e0uxb7ecuxb2dd-uxd559uxc2b5-uxacfcuxc815-10uxbd84}}

기계학습이 실제로 어떻게 작동하는지 4단계 과정을 알아봅시다.

\textbf{{[}그림 04{]} 머신러닝 학습 과정}

\begin{verbatim}
데이터 수집 → 학습 → 평가 → 예측
\end{verbatim}

\hypertarget{uxb2e8uxacc4-uxb370uxc774uxd130-uxc218uxc9d1-data-collection}{%
\subsection{1단계: 데이터 수집 (Data
Collection)}\label{uxb2e8uxacc4-uxb370uxc774uxd130-uxc218uxc9d1-data-collection}}

\begin{quote}
\textbf{``좋은 데이터 = 좋은 AI''}
\end{quote}

\hypertarget{uxb370uxc774uxd130uxb780}{%
\subsubsection{데이터란?}\label{uxb370uxc774uxd130uxb780}}

기계학습에서 데이터는 ``학습 재료''예요. 음식을 만들 때 재료가 좋아야
맛있는 요리가 나오듯, 좋은 데이터가 있어야 좋은 AI가 만들어집니다.

\hypertarget{uxb370uxc774uxd130-uxc218uxc9d1-uxc608uxc2dc}{%
\subsubsection{데이터 수집
예시}\label{uxb370uxc774uxd130-uxc218uxc9d1-uxc608uxc2dc}}

\textbf{고양이 vs 개 분류 모델 만들기}

\begin{enumerate}
\def\labelenumi{\arabic{enumi}.}
\tightlist
\item
  고양이 사진 1,000장 수집
\item
  개 사진 1,000장 수집
\item
  각 사진에 라벨 붙이기:

  \begin{itemize}
  \tightlist
  \item
    cat\_001.jpg → ``고양이''
  \item
    dog\_001.jpg → ``개''
  \end{itemize}
\end{enumerate}

\hypertarget{uxc88buxc740-uxb370uxc774uxd130uxc758-uxc870uxac74}{%
\subsubsection{좋은 데이터의
조건}\label{uxc88buxc740-uxb370uxc774uxd130uxc758-uxc870uxac74}}

\hypertarget{uxcda9uxbd84uxd55c-uxc591}{%
\paragraph{\emoji{check-mark-button} 1. 충분한 양}\label{uxcda9uxbd84uxd55c-uxc591}}

\begin{itemize}
\tightlist
\item
  적은 데이터: 100장 → 성능 낮음
\item
  많은 데이터: 10,000장 → 성능 높음
\item
  일반적으로 수백\textasciitilde 수천 개 필요
\end{itemize}

\hypertarget{uxb192uxc740-uxd488uxc9c8}{%
\paragraph{\emoji{check-mark-button} 2. 높은 품질}\label{uxb192uxc740-uxd488uxc9c8}}

\begin{itemize}
\tightlist
\item
  정확한 라벨 (고양이를 개라고 잘못 표시 X)
\item
  선명한 이미지 (흐릿한 사진 X)
\item
  오류가 없는 데이터
\end{itemize}

\hypertarget{uxb2e4uxc591uxc131}{%
\paragraph{\emoji{check-mark-button} 3. 다양성}\label{uxb2e4uxc591uxc131}}

\begin{itemize}
\tightlist
\item
  다양한 품종 포함
\item
  다양한 각도의 사진
\item
  다양한 배경, 조명
\end{itemize}

\hypertarget{uxd3b8uxd5a5uxb41c-uxb370uxc774uxd130-uxc8fcuxc758}{%
\paragraph{\emoji{cross-mark} 편향된 데이터
주의}\label{uxd3b8uxd5a5uxb41c-uxb370uxc774uxd130-uxc8fcuxc758}}

\textbf{나쁜 예시}: - 고양이는 실내 사진만 - 개는 야외 사진만 - → AI가
``실내 = 고양이''로 잘못 학습

\textbf{좋은 예시}: - 고양이: 실내 50\% + 야외 50\% - 개: 실내 50\% +
야외 50\%

\hypertarget{uxc5d4uxd2b8uxb9ac-uxc5f0uxacb0}{%
\subsubsection{엔트리 연결}\label{uxc5d4uxd2b8uxb9ac-uxc5f0uxacb0}}

엔트리 AI 모델에서: 1. 웹캠으로 직접 사진 촬영 (데이터 수집) 2. ``클래스
1'', ``클래스 2''로 분류 (라벨링) 3. 각 클래스당 20\textasciitilde30장
권장

\hypertarget{uxb2e8uxacc4-uxd559uxc2b5-training}{%
\subsection{2단계: 학습
(Training)}\label{uxb2e8uxacc4-uxd559uxc2b5-training}}

\begin{quote}
\textbf{컴퓨터가 데이터에서 패턴을 찾는 과정}
\end{quote}

\hypertarget{uxd559uxc2b5uxc774uxb780}{%
\subsubsection{학습이란?}\label{uxd559uxc2b5uxc774uxb780}}

컴퓨터가 수많은 예시를 보며 ``고양이는 이런 특징이 있구나'', ``개는 저런
특징이 있구나''를 스스로 발견하는 과정이에요.

\hypertarget{uxd559uxc2b5-uxacfcuxc815}{%
\subsubsection{학습 과정}\label{uxd559uxc2b5-uxacfcuxc815}}

\begin{enumerate}
\def\labelenumi{\arabic{enumi}.}
\tightlist
\item
  \textbf{초기 상태}: 아무것도 모름

  \begin{itemize}
  \tightlist
  \item
    고양이인지 개인지 랜덤으로 맞춤
  \item
    정확도 50\% (동전 던지기 수준)
  \end{itemize}
\item
  \textbf{학습 시작}: 데이터를 반복해서 봄

  \begin{itemize}
  \tightlist
  \item
    1회 (Epoch 1): ``아, 귀 모양이 중요하구나''
  \item
    2회 (Epoch 2): ``수염도 중요하네''
  \item
    3회 (Epoch 3): ``눈 모양도 봐야겠다''
  \end{itemize}
\item
  \textbf{가중치 조정}: 중요한 특징에 더 큰 점수 부여

  \begin{itemize}
  \tightlist
  \item
    귀 모양: 가중치 높임
  \item
    털 색깔: 가중치 낮춤 (고양이/개 모두 다양)
  \end{itemize}
\item
  \textbf{반복 학습}:

  \begin{itemize}
  \tightlist
  \item
    10회, 20회, 50회\ldots{} 반복
  \item
    점점 정확도 향상
  \end{itemize}
\end{enumerate}

\hypertarget{uxd559uxc2b5-uxc2dcuxac04}{%
\subsubsection{학습 시간}\label{uxd559uxc2b5-uxc2dcuxac04}}

\begin{longtable}{p{4.5cm}p{4.5cm}p{5.0cm}}
\toprule\noalign{}
데이터 크기 & 모델 복잡도 & 학습 시간 \\
\midrule\noalign{}
\endhead
\bottomrule\noalign{}
\endlastfoot
작음 & 간단 & 수 초 \textasciitilde{} 수 분 \\
보통 & 보통 & 수 시간 \textasciitilde{} 수 일 \\
큼 & 복잡 (딥러닝) & 수 주 \textasciitilde{} 수 개월 \\
\end{longtable}

ChatGPT 같은 대형 모델은 수천 개의 GPU로 수개월간 학습했어요!

\hypertarget{uxc5d4uxd2b8uxb9ac-uxc5f0uxacb0-1}{%
\subsubsection{엔트리 연결}\label{uxc5d4uxd2b8uxb9ac-uxc5f0uxacb0-1}}

엔트리 AI 모델: 1. ``모델 학습하기'' 버튼 클릭 2. 진행 바가 100\%까지
진행 3. 수십 초 \textasciitilde{} 수 분 소요 (간단한 모델이라 빠름)

\hypertarget{uxb2e8uxacc4-uxd3c9uxac00-testing}{%
\subsection{3단계: 평가
(Testing)}\label{uxb2e8uxacc4-uxd3c9uxac00-testing}}

\begin{quote}
\textbf{학습한 모델이 얼마나 정확한지 확인}
\end{quote}

\hypertarget{uxd3c9uxac00uxb780}{%
\subsubsection{평가란?}\label{uxd3c9uxac00uxb780}}

시험을 본다고 생각하면 돼요. 학습할 때 보지 못한 \textbf{새로운
데이터}로 테스트합니다.

\hypertarget{uxc65c-uxc0c8uxb85cuxc6b4-uxb370uxc774uxd130uxb85c-uxd3c9uxac00uxd560uxae4c}{%
\subsubsection{왜 새로운 데이터로
평가할까?}\label{uxc65c-uxc0c8uxb85cuxc6b4-uxb370uxc774uxd130uxb85c-uxd3c9uxac00uxd560uxae4c}}

\textbf{나쁜 평가 방법}: 학습 데이터로 평가 - 학생이 시험에 나온 문제를
미리 외워버림 - 100점을 받아도 진짜 실력인지 알 수 없음

\textbf{좋은 평가 방법}: 새로운 데이터로 평가 - 처음 보는 문제로 시험 -
진짜 실력을 측정 가능

\hypertarget{uxb370uxc774uxd130-uxbd84uxd560}{%
\subsubsection{데이터 분할}\label{uxb370uxc774uxd130-uxbd84uxd560}}

일반적으로 데이터를 다음과 같이 나눕니다:

\begin{verbatim}
전체 데이터 1,000장
 ├─ 학습 데이터 (Training): 700장 (70%)
 ├─ 검증 데이터 (Validation): 150장 (15%)
 └─ 테스트 데이터 (Testing): 150장 (15%)
\end{verbatim}

\begin{itemize}
\tightlist
\item
  \textbf{학습 데이터}: 모델 학습에 사용
\item
  \textbf{검증 데이터}: 학습 중간에 성능 확인
\item
  \textbf{테스트 데이터}: 최종 성능 평가
\end{itemize}

\hypertarget{uxd3c9uxac00-uxc9c0uxd45c-uxc815uxd655uxb3c4-accuracy}{%
\subsubsection{평가 지표: 정확도
(Accuracy)}\label{uxd3c9uxac00-uxc9c0uxd45c-uxc815uxd655uxb3c4-accuracy}}

\textbf{정확도 = (맞춘 개수 / 전체 개수) × 100\%}

예시: - 테스트 데이터 100장 - 90장 맞춤, 10장 틀림 - 정확도 = 90\%

\hypertarget{uxc131uxb2a5-uxac1cuxc120}{%
\subsubsection{성능 개선}\label{uxc131uxb2a5-uxac1cuxc120}}

정확도가 낮으면? 1. 데이터 더 수집 2. 모델 구조 변경 3. 학습 시간 늘리기
4. 다시 학습 (2단계로 돌아가기)

\hypertarget{uxc5d4uxd2b8uxb9ac-uxc5f0uxacb0-2}{%
\subsubsection{엔트리 연결}\label{uxc5d4uxd2b8uxb9ac-uxc5f0uxacb0-2}}

엔트리 AI 모델: 1. 학습 완료 후 정확도 표시 (예: 92\%) 2. 웹캠으로
실시간 테스트 3. 잘못 분류하면 데이터 추가 후 재학습

\hypertarget{uxb2e8uxacc4-uxc608uxce21-prediction}{%
\subsection{4단계: 예측
(Prediction)}\label{uxb2e8uxacc4-uxc608uxce21-prediction}}

\begin{quote}
\textbf{학습한 모델을 실전에 적용}
\end{quote}

\hypertarget{uxc608uxce21uxc774uxb780}{%
\subsubsection{예측이란?}\label{uxc608uxce21uxc774uxb780}}

드디어 완성된 AI 모델을 실제로 사용하는 단계예요!

\hypertarget{uxc608uxce21-uxacfcuxc815}{%
\subsubsection{예측 과정}\label{uxc608uxce21-uxacfcuxc815}}

\begin{enumerate}
\def\labelenumi{\arabic{enumi}.}
\tightlist
\item
  \textbf{새로운 입력}: 처음 보는 고양이 사진
\item
  \textbf{모델 적용}: 학습한 패턴으로 분석
\item
  \textbf{결과 출력}: ``고양이입니다'' (확률 95\%)
\end{enumerate}

\hypertarget{uxc2e4uxc804-uxbc30uxd3ec}{%
\subsubsection{실전 배포}\label{uxc2e4uxc804-uxbc30uxd3ec}}

완성된 모델을 다양한 곳에 활용:

\begin{longtable}{p{4cm}p{10cm}}
\toprule\noalign{}
분야 & 활용 예시 \\
\midrule\noalign{}
\endhead
\bottomrule\noalign{}
\endlastfoot
\emoji{mobile-phone} \textbf{앱} & 구글 포토 (사진 자동 분류) \\
\emoji{globe-with-meridians} \textbf{웹사이트} & Gmail (스팸 필터) \\
\emoji{robot} \textbf{로봇} & 로봇 청소기 (장애물 인식) \\
\emoji{automobile} \textbf{자동차} & 테슬라 (자율주행) \\
\emoji{hospital} \textbf{병원} & AI 영상 진단 시스템 \\
\end{longtable}

\hypertarget{uxc9c0uxc18duxc801uxc778-uxbaa8uxb2c8uxd130uxb9c1}{%
\subsubsection{지속적인
모니터링}\label{uxc9c0uxc18duxc801uxc778-uxbaa8uxb2c8uxd130uxb9c1}}

배포 후에도 계속 관리가 필요해요:

\begin{enumerate}
\def\labelenumi{\arabic{enumi}.}
\tightlist
\item
  \textbf{성능 모니터링}: 정확도가 떨어지지 않는지 확인
\item
  \textbf{새 데이터 수집}: 환경이 변하면 새 데이터 추가
\item
  \textbf{재학습}: 주기적으로 모델 업데이트
\end{enumerate}

예시: Gmail 스팸 필터 - 매일 새로운 스팸 기법 등장 - 지속적으로 데이터
수집 및 재학습 - 모델 계속 개선

\hypertarget{uxc5d4uxd2b8uxb9ac-uxc5f0uxacb0-3}{%
\subsubsection{엔트리 연결}\label{uxc5d4uxd2b8uxb9ac-uxc5f0uxacb0-3}}

엔트리 AI 모델: 1. 웹캠으로 새 이미지 입력 2. 실시간으로 분류 결과 표시
3. 블록 코딩과 연동하여 작품 제작 - ``만약 고양이면 → 야옹 소리'' -
``만약 개면 → 멍멍 소리''

\hypertarget{uxb2e8uxacc4-uxacfcuxc815-uxc694uxc57d}{%
\subsection{4단계 과정
요약}\label{uxb2e8uxacc4-uxacfcuxc815-uxc694uxc57d}}

\begin{longtable}{p{2.0cm}p{4.0cm}p{4.0cm}p{4.0cm}}
\toprule\noalign{}
단계 & 이름 & 설명 & 엔트리 \\
\midrule\noalign{}
\endhead
\bottomrule\noalign{}
\endlastfoot
1 & 데이터 수집 & 학습할 재료 모으기 & 웹캠으로 사진 촬영 \\
2 & 학습 & 패턴 발견 & ``모델 학습하기'' 버튼 \\
3 & 평가 & 정확도 확인 & 정확도 \% 표시 \\
4 & 예측 & 실전 적용 & 새 이미지 분류 \\
\end{longtable}

\begin{quote}
\emoji{light-bulb} \textbf{핵심 메시지}: 기계학습은 ``데이터 수집 → 학습 → 평가 →
예측''의 4단계로 이루어집니다. 3주차 엔트리 실습에서 이 전 과정을 직접
체험하게 됩니다!
\end{quote}

\hypertarget{uxc815uxb9ac-5uxbd84}{%
\section{\emoji{memo} 정리 (5분)}\label{uxc815uxb9ac-5uxbd84}}

\hypertarget{uxd575uxc2ec-uxac1cuxb150-uxc694uxc57d}{%
\subsection{핵심 개념
요약}\label{uxd575uxc2ec-uxac1cuxb150-uxc694uxc57d}}

오늘 배운 내용을 4가지로 정리해 봅시다.

\hypertarget{ai-ml-dl-uxd3ecuxd568-uxad00uxacc4}{%
\subsubsection{\textbf{①} AI ⊃ ML ⊃ DL 포함
관계}\label{ai-ml-dl-uxd3ecuxd568-uxad00uxacc4}}

\begin{verbatim}
AI (인공지능) - 가장 큰 개념
 ⊃ ML (기계학습) - 데이터로 학습하는 AI
    ⊃ DL (딥러닝) - 신경망이 깊은 ML
\end{verbatim}

\begin{itemize}
\tightlist
\item
  AI는 인간의 지능을 모방하는 모든 기술
\item
  기계학습은 AI를 구현하는 대표적 방법
\item
  딥러닝은 기계학습의 한 종류
\end{itemize}

\hypertarget{uxae30uxacc4uxd559uxc2b5uxc758-uxc815uxc758}{%
\subsubsection{\textbf{②} 기계학습의
정의}\label{uxae30uxacc4uxd559uxc2b5uxc758-uxc815uxc758}}

\begin{quote}
\textbf{기계학습}: 데이터로부터 패턴을 스스로 학습하는 AI 기술
\end{quote}

\textbf{전통적 프로그래밍 vs 기계학습} - 전통적: 규칙 작성 → 결과 -
기계학습: 데이터 + 결과 → 규칙 발견

\hypertarget{uxac00uxc9c0-uxd559uxc2b5-uxc720uxd615}{%
\subsubsection{\textbf{③} 3가지 학습
유형}\label{uxac00uxc9c0-uxd559uxc2b5-uxc720uxd615}}

\begin{longtable}{p{2.5cm}p{3.5cm}p{4.0cm}p{4.0cm}}
\toprule\noalign{}
\textbf{유형} & \textbf{특징} & \textbf{비유} & \textbf{대표 예시} \\
\midrule\noalign{}
\endhead
\bottomrule\noalign{}
\endlastfoot
\textbf{지도학습} & 정답 있는 데이터로 학습 & 시험 공부 & 스팸 필터,
질병 진단 \\
\textbf{비지도학습} & 정답 없이 패턴 발견 & 그룹 나누기 & 고객 세분화 \\
\textbf{강화학습} & 시행착오로 보상 최대화 & 게임 연습 & AlphaGo,
자율주행 \\
\end{longtable}

\textbf{\emoji{light-bulb} 초등 AI 교육에서는 주로 지도학습을 다룹니다!}

\hypertarget{uxba38uxc2e0uxb7ecuxb2dd-4uxb2e8uxacc4-uxacfcuxc815}{%
\subsubsection{\textbf{④} 머신러닝 4단계
과정}\label{uxba38uxc2e0uxb7ecuxb2dd-4uxb2e8uxacc4-uxacfcuxc815}}

\begin{verbatim}
1. 데이터 수집 → 2. 학습 → 3. 평가 → 4. 예측
\end{verbatim}

\begin{itemize}
\tightlist
\item
  \textbf{1단계}: 좋은 데이터 모으기 (양, 질, 다양성)
\item
  \textbf{2단계}: 패턴 발견 (반복 학습, 가중치 조정)
\item
  \textbf{3단계}: 성능 확인 (새 데이터로 테스트)
\item
  \textbf{4단계}: 실전 적용 (새로운 입력 예측)
\end{itemize}

\hypertarget{uxc5d4uxd2b8uxb9ac-ai-uxbaa8uxb378-uxc5f0uxacb0-1}{%
\subsection{엔트리 AI 모델
연결}\label{uxc5d4uxd2b8uxb9ac-ai-uxbaa8uxb378-uxc5f0uxacb0-1}}

오늘 배운 이론이 3주차 엔트리 실습으로 연결됩니다!

\textbf{엔트리 이미지 분류 모델}: - \textbf{학습 유형}: 지도학습 -
\textbf{1단계}: 웹캠으로 고양이/개 사진 촬영 - \textbf{2단계}: ``모델
학습하기'' 버튼 클릭 - \textbf{3단계}: 정확도 확인 (예: 92\%) -
\textbf{4단계}: 새 이미지 실시간 분류

→ \textbf{이 모든 것이 바로 지도학습입니다!}

\hypertarget{uxb2e4uxc74c-uxc8fc-uxd034uxc988-uxc548uxb0b4}{%
\subsection{다음 주 퀴즈
안내}\label{uxb2e4uxc74c-uxc8fc-uxd034uxc988-uxc548uxb0b4}}

\textbf{2주차 1차시 시작 시 1주차 전체 퀴즈 실시}

\begin{longtable}{p{4cm}p{10cm}}
\toprule\noalign{}
항목 & 내용 \\
\midrule\noalign{}
\endhead
\bottomrule\noalign{}
\endlastfoot
\textbf{시간} & 10분 \\
\textbf{문항} & 5지선다형 10문항 + 단답형 5문항 \\
\textbf{배점} & 10점 (5지선다 0.5점 × 10 + 단답형 1점 × 5) \\
\textbf{출제 범위} & 1주차 1차시 + 1주차 2차시 전체 \\
\end{longtable}

\textbf{복습 포인트}: - 1차시: AI 정의, 약한/강한 AI, 기능별 분류, AI
봄과 겨울 - 2차시: AI/ML/DL 포함 관계, 3가지 학습 유형, 머신러닝 4단계

\textbf{복습 권장!}

\hypertarget{uxb2e4uxc74c-uxcc28uxc2dc-uxc608uxace0}{%
\subsection{다음 차시
예고}\label{uxb2e4uxc74c-uxcc28uxc2dc-uxc608uxace0}}

\textbf{2주차 1차시: 데이터와 딥러닝}

\begin{itemize}
\tightlist
\item
  데이터의 역할과 중요성
\item
  데이터 편향 문제
\item
  신경망과 딥러닝의 기초
\item
  AI 시대의 데이터 윤리
\end{itemize}

\textbf{준비물}: 없음 (이론 강의)

\hypertarget{uxcc38uxace0uxc790uxb8cc}{%
\subsection{참고자료}\label{uxcc38uxace0uxc790uxb8cc}}

\begin{longtable}{p{2.5cm}p{3.5cm}p{4.0cm}p{4.0cm}}
\toprule\noalign{}
\textbf{\#} & \textbf{자료명} & \textbf{유형} & \textbf{비고} \\
\midrule\noalign{}
\endhead
\bottomrule\noalign{}
\endlastfoot
1 & AI디지털리터러시\_강의계획서\_v1.2.0.0 & 프로젝트 문서 & 1주차 2차시
수업 계획 \\
2 & AI디지털리터러시\_1주차2차시\_강의개요\_v1.2.0.0 & 프로젝트 문서 &
수업 진행 가이드 \\
3 & 2022 개정 교육과정 (교육부 고시 제2022-33호) & 교육부 문서 &
{[}6실05-05{]} 성취기준 \\
4 & AI 기초 (한국정보교육학회 편) & 도서 & 기계학습 원리 참고 \\
5 & 교수 제공 AI/ML 시뮬레이터 & 교수 제공 & 지도/비지도학습 시연 \\
6 & 강화학습 테트리스 시뮬레이터 & 교수 제공 & 강화학습 개념 시연 \\
7 & AI과수원 애니메이션 & 교수 제공 & 기계학습 과정 시각화 \\
8 & 엔트리 (playentry.org) & 웹사이트 & 이미지 분류 모델 예시 \\
9 & 문서작성지침서\_v1.7.2.0 & 프로젝트 문서 & 강의자료 작성 절차 \\
10 & 버전관리지침\_v1.3.0 & 프로젝트 문서 & 문서 버전 체계 \\
\end{longtable}

\hypertarget{uxc0dduxc131uxc218uxc815-uxc774uxb825}{%
\subsection{생성/수정
이력}\label{uxc0dduxc131uxc218uxc815-uxc774uxb825}}

\begin{longtable}{p{1.5cm}p{2.0cm}p{2.5cm}p{2.5cm}p{2.5cm}p{3.0cm}}
\toprule\noalign{}
\textbf{버전} & \textbf{날짜} & \textbf{시간} & \textbf{변경 수준} & \textbf{변경 내용} & \textbf{작성자} \\
\midrule\noalign{}
\endhead
\bottomrule\noalign{}
\endlastfoot
v1.0.0.0 & 2026-02-08 & 16:30 & 최초 & 1주차 2차시 강의노트 신규 작성.
1주차 1차시 구조(이모지, 상세 설명, 표 활용) 참조. AI/ML/DL 포함관계,
3가지 학습 유형(지도/비지도/강화) 상세 설명, 머신러닝 4단계 프로세스
정리, 엔트리 AI 모델 연결 강조, 일상 예시 풍부하게 제공 & Claude \\
v1.2.0.0 & 2026-02-08 & 19:30 & 마이너 & 비지도학습 ``정답이 없다''의
의미 명확화. ① ``정답 없이'' → ``라벨(정답 표시) 없이''로 수정 ②
지도학습과의 데이터 구조 비교 추가 ③ 라벨이 없는 3가지 이유 상세 설명
(예시 확대) ④ 군집화 3단계 과정 구조화 ⑤ 구체적 예시 3개 확대 및 재구성
⑥ 초등교육 비유 2가지 추가 ⑦ 오개념 대응 전략 표 신설 ⑧ 참고자료 추가
(테트리스 시뮬레이터, AI과수원) & Claude \\
v1.2.1.0 & 2026-02-08 & 20:45 & 리비전 & 전문가 검토 의견 반영. ①
회귀(Regression) 용어 구체화: ``시험 점수 맞히기 vs 합격/불합격'' 비유
추가, 과일 무게 예측 비유 추가 ② 특징(Feature) 개념 확장: 고양이/개 특징
데이터 표 추가, AI가 발견하는 패턴 명시, 특징 추출 과정 상세화 ③ 초보자
및 예비 교사 관점 접근성 강화 & Claude \\
\textbf{v1.3.0.0} & \textbf{2026-02-08} & \textbf{21:45} &
\textbf{마이너} & \textbf{완전한 개선 반영 (전문가 검토 의견 3가지 모두
반영). ① 회귀 용어 완전 구체화: ``얼마나?'' 질문 추가, 회귀 vs 분류
비교표 신설, 교수 멘트 추가, 4개 예시 확대 ② 강화학습 시각화 완성: 그림
19(강아지 훈련) 추가, 그림 20(게임 점수) 추가, 교수 전략 상세화,
초등학생 설명 추가, 오개념 방지 표 신설 ③ 교육적 접근성 극대화: 시각
자료 2개 추가, 일상 비유 확대, 단계별 교수 멘트 보강} &
\textbf{Claude} \\
\end{longtable}

\begin{center}\rule{0.5\linewidth}{0.5pt}\end{center}

\emph{AI디지털리터러시 1주차 2차시 강의노트 --- 기계학습의 원리}

\end{document}
